%% kapa: exam=1 round=33 year=2022 period=2 form=A subject=accounting subject_ko=회계학 seq=38 number=78
%% source: Q-Net 감정평가사 기출문제 (cst00309) artlSeq=5209997 file=raw/1008/5209997/2022년도 제33회 감정평가사 1차시험 2교시(감정평가 관계법규, 회계학).pdf
%% conversion: kordoc 4.24.51 -> tools/gen1.py (auto); review=proofread-sonnet; flags=none
\begin{question}{78}
다음은 제품 A를 생산·판매하는 (주)감평의 당기 전부원가 손익계산서와 공헌이익 손익계산서이다.
\begin{framedtext}
{\small\begin{tabular}{@{}lr|lr@{}}
\multicolumn{2}{c|}{\uline{전부원가 손익계산서}} & \multicolumn{2}{c}{\uline{공헌이익 손익계산서}} \\
매출액 & ￦1,000,000 & 매출액 & ￦1,000,000 \\
매출원가 & \uline{650,000} & 변동원가 & \uline{520,000} \\
매출총이익 & 350,000 & 공헌이익 & 480,000 \\
판매관리비 & \uline{200,000} & 고정원가 & \uline{400,000} \\
영업이익 & \uline{150,000} & 영업이익 & \uline{80,000} \\
\end{tabular}\par}
\end{framedtext}
제품의 단위당 판매가격 ￦1,000, 총고정판매관리비가 ￦50,000일 때 전부원가계산에 의한 기말제품재고는? (단, 기초 및 기말 재공품, 기초제품은 없다.)
\begin{choices}
  \item ￦85,000
  \item ￦106,250
  \item ￦162,500
  \item ￦170,000
  \item ￦212,500
\end{choices}
\end{question}
\answer{3}
